\documentclass[11pt]{article}

\usepackage[T1]{fontenc}
\usepackage[utf8]{inputenc}
\usepackage{lmodern}
\usepackage{amsmath,amssymb,amsthm}
\usepackage{booktabs}
\usepackage{array}
\usepackage{xcolor}
\usepackage{geometry}
\usepackage{enumitem}
\usepackage{hyperref}

\geometry{margin=1in}
\hypersetup{
  colorlinks=true,
  linkcolor=blue!45!black,
  citecolor=blue!45!black,
  urlcolor=blue!55!black,
  pdftitle={Certified Bounded Revision},
  pdfsubject={Safety and finite-sample risk control for autonomous long-form generation}
}
\setlist{leftmargin=*,itemsep=0.25em,topsep=0.4em}
\setlength{\parskip}{0.35em}
\setlength{\parindent}{1.2em}
\raggedbottom

\pagestyle{plain}

\newtheorem{definition}{Definition}[section]
\newtheorem{theorem}[definition]{Theorem}
\newtheorem{proposition}[definition]{Proposition}
\newtheorem{lemma}[definition]{Lemma}
\newtheorem{corollary}[definition]{Corollary}
\newtheorem{assumption}[definition]{Assumption}
\newtheorem{remark}[definition]{Remark}

\newcommand{\code}[1]{\texttt{#1}}
\newcommand{\Prob}{\mathbb{P}}
\newcommand{\E}{\mathbb{E}}
\newcommand{\N}{\mathbb{N}}
\newcommand{\ind}{\mathbf{1}}
\newcommand{\sha}{\operatorname{SHA256}}
\newcommand{\BetaInv}{\operatorname{Beta}^{-1}}

\title{\textbf{Certified Bounded Revision}\\[0.35em]
\large Deterministic Safety and Finite-Sample Risk Control\\
for Autonomous Long-Form Generation}
\author{Technical report prepared for the LitHarness project}
\date{17 August 2026}

\begin{document}
\maketitle

\begin{abstract}
Autonomous long-form generation has an asymmetric evaluation problem. It is easy to
verify that a manuscript exists, that a patch is syntactically bounded, or that a ledger
balances. It is substantially harder to verify dramatic function, voice, escalation, and
payoff. Treating a noisy quality proxy as an objective creates an optimization channel;
treating it as a blocking test creates an evidence problem. This paper develops a safety
case for an autonomous manuscript-revision system that keeps these problems separate.

We formalize a manuscript as a content-addressed ordered state and a repair as an
interval-certified partial transformation. We prove four deterministic properties:
localized noninterference outside licensed spans, rejection of stale patches, atomic
auditability of accepted revisions, and finite repair cascades. For the concrete
LitHarness bounds, repair depth $D=3$ and propagation fan-out $F=12$ imply that one root
evaluation can create at most 55 evaluation and repair jobs, before idempotent
deduplication. We then model a craft gate as selective prediction with abstention and show
that the current point-estimate rule, observed precision at least 0.80 with at least 17
flags, does not establish a 0.80 precision guarantee. An executable counterexample with
14 correct flags out of 17 is promotable while its exact two-sided 95\% Clopper-Pearson
lower limit is 0.566. We propose a schema-complete promotion certificate that stores
integer confusion counts, separates threshold selection from safety testing, and requires
an exact lower confidence bound, with family-wise correction when several thresholds are
tested.

The result is deliberately narrower than a proof of literary quality. It is a proof that
uncertified generators cannot silently rewrite unrelated prose, cannot close findings on
detector failure, cannot recurse without bound, and cannot turn a quality hypothesis into
a blocking policy without reconstructible statistical evidence. That narrower claim is
both enforceable and useful.
\end{abstract}

\noindent\textbf{Keywords:} autonomous agents; long-form generation; runtime assurance;
selective prediction; exact binomial confidence bounds; bounded repair; content-addressed
state; Goodhart effects.

\section{Introduction}

Long-form generation is not merely repeated next-token prediction. A book-production
system must preserve commitments across scenes, maintain a changing world state, apply
director constraints, repair local defects, and decide when generated text becomes canon.
The decision boundary is the difficult part. A language model can propose a paragraph, but
its proposal is not evidence that the paragraph is correct, safe to splice into the book,
or worth keeping.

The usual response is to add a critic. For open-ended prose this substitutes one uncertain
model for another. LLM judges are known to exhibit position, verbosity, and self-enhancement
biases \cite{zheng2023judge}, and intrinsic self-correction can degrade reasoning without
external feedback \cite{huang2024selfcorrect}. The LitHarness project supplies a sharper
empirical warning. Twenty static quality proxies were rejected under controls; candidate
signals measured era, annotation density, chapter count, compression-window mechanics, or
sampling coverage rather than the intended defect \cite{litharnessbrief}. The lesson is not
that measurement is futile. It is that a measurement must earn the particular decision it
is allowed to make.

This paper asks a different question:

\begin{quote}
What can be certified about an autonomous revision process even when prose quality itself
cannot yet be certified?
\end{quote}

Our answer is a two-layer architecture.

\begin{enumerate}
  \item A small deterministic safety kernel constrains where and how generated text may
  change the manuscript. Its properties are universal over generator behavior.
  \item Statistical craft gates are optional selective decisions. They may block only when
  the evidence record supports an exact finite-sample risk statement at the same unit of
  analysis as the blocked decision.
\end{enumerate}

This resembles runtime-enforcement work, in which an untrusted action stream is mediated by
a trusted monitor \cite{ligatti2005edit}, and Seldonian design, in which the output may be
``no solution found'' unless a high-probability behavioral constraint is established
\cite{thomas2019seldonian}. The prose setting changes the specification. We do not have a
formal predicate for ``good scene.'' We do have formal predicates for scope preservation,
version freshness, transaction completeness, bounded scheduling, and the provenance of a
statistical claim.

\subsection{Contributions}

The paper makes five contributions.

\begin{enumerate}
  \item It gives a formal model of content-addressed manuscript revision and a repair
  certificate whose accepted effect is confined to cited intervals.
  \item It proves localized noninterference, stale-write rejection, and atomic auditability
  for the implemented repair path.
  \item It derives a finite upper bound for detect-repair-propagate cascades. The bound is
  expressed by a second-order recurrence and instantiated from the repository constants.
  \item It identifies a finite-sample gap in the current craft-promotion rule and provides an
  executable counterexample. The proposed correction uses exact binomial lower bounds and
  multiplicity-aware candidate selection.
  \item It formalizes an optimization firewall: a craft failure parks the unit instead of
  feeding a proxy into repeated generation, while deterministic defects may license bounded
  retries or located repair.
\end{enumerate}

\subsection{Non-claims}

The following distinctions are load-bearing.

\begin{itemize}
  \item Preservation is not quality. Byte-identical untouched prose may still be poor.
  \item Detector completion is not detector soundness or completeness. A complete run can
  miss a defect.
  \item A finite cascade is not a successful cascade. Termination can end in a parked job.
  \item Precision among flagged units is not recall, coverage, utility, or reader value.
  \item A passed deterministic gate means ``no encoded violation was found,'' not ``the
  scene is good.''
\end{itemize}

\section{System Model and Threat Model}

\subsection{Manuscripts, versions, and heads}

Let $\Sigma^*$ be the set of finite byte strings after a declared canonicalization function.
A manuscript revision is
\[
  M = (B, R, p, \mathcal{V}),
\]
where $B$ is a book identifier, $R$ is a branch identifier, $p$ is an optional parent
revision identifier, and $\mathcal{V}=(v_1,\ldots,v_n)$ is an ordered finite tuple of nodes.
A node is
\[
  v_i=(\ell_i,\pi_i,k_i,c_i,h_i,z_i),
\]
with stable logical identifier $\ell_i$, position key $\pi_i$, node kind $k_i$, content
$c_i\in\Sigma^*\cup\{\bot\}$, content hash $h_i=\sha(c_i)$ when content exists, and lock
state $z_i$. The node version identifier is a digest of its serialized fields. The revision
identifier is a digest of the ordered node versions and branch metadata.

The mutable operational state contains a single head pointer $H(B,R)$ to an immutable
revision. Moving the head forward creates history; it does not overwrite an earlier
revision. This is compatible with optimistic concurrency control: work proceeds against a
read version and validates that version before commit \cite{kung1981optimistic}.

\subsection{Generated components and trusted components}

The generator, critic, and information extractor are \emph{untrusted} for safety purposes.
They may return malformed output, choose an incorrect span, cite stale text, propose a
whole-scene rewrite, or systematically optimize a shallow proxy. Provider failures and
timeouts are also permitted.

The trusted computing base for the claims in this paper consists of:

\begin{itemize}
  \item canonical serialization and cryptographic hashing;
  \item interval and policy checks in the patch applicator;
  \item the pure decision mapping from gate results to workflow outcomes;
  \item transaction semantics of the SQLite store;
  \item content-derived job identity and the explicit depth and fan-out counters; and
  \item the statistical certificate checker proposed in Section~\ref{sec:statistics}.
\end{itemize}

Cryptographic collision resistance is assumed. We do not claim protection against a hostile
process that can arbitrarily modify the trusted code, database engine, or host filesystem.

\subsection{Adversarial and accidental failure modes}

The threat model includes five classes.

\begin{enumerate}
  \item \textbf{Spatial overreach:} a local complaint licenses a global rewrite.
  \item \textbf{Temporal staleness:} a valid patch is applied after the target has changed.
  \item \textbf{False completion:} an evaluator error is interpreted as zero findings.
  \item \textbf{Unbounded work:} repair, re-evaluation, and propagation create a cycle or
  combinatorial queue explosion.
  \item \textbf{Epistemic overreach:} a population statistic or point estimate is recorded
  as evidence of scene quality, or a critic is optimized against until it passes.
\end{enumerate}

The first four admit deterministic controls. The fifth requires a statistical protocol and
an explicit refusal to optimize the measuring instrument.

\section{Interval-Certified Repair}

\subsection{Repair certificates}

\begin{definition}[Repair certificate]
For target node $v=(\ell,\pi,k,c,h,z)$, a repair certificate is
\[
 P=(\ell,\nu,h,\{(a_j,b_j,s_j,g_j)\}_{j=1}^m,f),
\]
where $\nu$ is the expected node-version identifier, $h$ is the expected base-content hash,
$[a_j,b_j)$ is a half-open interval in $c$, $s_j$ is replacement text, $g_j$ is an optional
hash of the cited substring, and $f$ is an optional finding identifier that licenses the
operation.
\end{definition}

The deterministic policy accepts $P$ only when all of the following hold.

\begin{enumerate}[label=(C\arabic*)]
  \item The target exists, has content, and is not content-locked.
  \item $\nu$ and $h$ equal the current node version and content hash.
  \item $0\le a_j\le b_j\le |c|$ for every $j$.
  \item If $g_j$ is present, then $g_j=\sha(c[a_j:b_j])$.
  \item Intervals are pairwise disjoint.
  \item A deletion is licensed when deletion licensing is required.
  \item The number of operations is at most $m_{\max}$ and
  $\sum_j(b_j-a_j)/|c|\le\rho$.
  \item The result-length ratio lies in $[\lambda_{\min},\lambda_{\max}]$.
\end{enumerate}

LitHarness currently uses $m_{\max}=32$, $\rho=0.5$,
$\lambda_{\min}=0.5$, and $\lambda_{\max}=2.0$. These are policy values, not empirical
claims about ideal prose revision.

Sort the intervals by start position and decompose the original content uniquely as
\[
 c=u_0\,x_1\,u_1\,x_2\cdots x_m\,u_m,
\]
where $x_j=c[a_j:b_j]$ and each $u_j$ is untouched. The patch result is
\[
 T_P(c)=u_0\,s_1\,u_1\,s_2\cdots s_m\,u_m.
\]

\subsection{Localized noninterference}

\begin{theorem}[Localized noninterference]
\label{thm:noninterference}
If a repair certificate $P$ is accepted, then every untouched segment $u_j$ occurs in
$T_P(c)$ byte-identically and in the same order. No accepted repair can insert, delete, or
substitute a byte outside the union of the cited intervals.
\end{theorem}

\begin{proof}
Conditions (C3) and (C5) make the sorted intervals a disjoint ordered family. Therefore the
decomposition $c=u_0x_1u_1\cdots x_mu_m$ exists and is unique. The splice construction copies
each $u_j$ directly from $c$ and substitutes only $x_j$ with $s_j$. The preservation verifier
independently walks the result using the cumulative replacement-length changes and compares
the result slice with each $u_j$. A mismatch rejects the patch. Hence acceptance implies
byte equality of every untouched segment, and their concatenation order is unchanged.
\end{proof}

\begin{remark}[Offsets versus identity]
Replacement lengths may differ, so an untouched suffix need not retain its numeric offset.
The invariant is segment identity and order under the induced offset map, not equality of
absolute offsets. This distinction avoids the common but false claim that a variable-length
patch leaves all later coordinates unchanged.
\end{remark}

\begin{corollary}[Bounded edit surface]
At most a fraction $\rho$ of the original node is licensed as a target, regardless of whether
replacement text has the same length. Thus a same-length whole-scene rewrite is rejected when
$\rho<1$.
\end{corollary}

The corollary matters because length checks alone do not distinguish a repair from a rewrite.
A generated scene of length $n$ can be replaced by unrelated text of the same length while
passing every global length-ratio test. Interval coverage closes that channel.

\subsection{Stale-patch rejection}

\begin{theorem}[Stale-write safety]
\label{thm:stale}
Suppose patch $P$ was constructed against node $v$ with version $\nu$ and content hash $h$.
If the current target node $v'$ has a different version identifier or content hash, then
$P$ is rejected before splicing.
\end{theorem}

\begin{proof}
Condition (C2) is necessary for acceptance. If either equality fails, the policy emits a
stale-base veto. The applicator returns without calling the splice operation. Therefore the
current content is unchanged by $P$.
\end{proof}

The version and full-content checks are complemented by per-span hashes. The full hash
protects the base object; the span hash catches a caller that presents inconsistent interval
evidence even when coordinates remain in range. Under collision resistance, an accepted
patch is bound to the exact bytes it claims to replace.

\subsection{Atomic auditability}

Let an accepted transition write the tuple
\[
 W=(M',H',S,E,D,J),
\]
where $M'$ is the immutable revision, $H'$ the branch head, $S$ extracted and retracted
state, $E$ events, $D$ the policy decision, and $J$ follow-on jobs.

\begin{theorem}[Atomic repair acceptance]
\label{thm:atomic}
Under atomic database transaction semantics, an accepted repair committed through the
LitHarness repair path has only two externally visible outcomes: none of $W$ is committed,
or all of $W$ is committed. In particular, the head cannot point to the accepted repair
without its policy decision and scheduled verification job.
\end{theorem}

\begin{proof}
The repair handler passes the new revision, head move, extracted state, retractions, events,
policy decision, verification job, and propagation jobs to one \code{commit\_revision}
invocation. That method encloses every corresponding write in one \code{BEGIN IMMEDIATE}
transaction and commits only after all writes and the head update succeed. On an exception,
the transaction rolls back. Atomicity therefore gives the stated all-or-none result.
\end{proof}

\begin{remark}
The theorem covers the implemented accepted-repair path. It is not a claim that every
possible future caller of the general storage method supplies a policy decision. The
artifact mapping in Section~\ref{sec:artifact} states the exact boundary.
\end{remark}

\section{Detect, Repair, Verify, and Propagate}

\subsection{Workflow semantics}

Let $E_d$ denote an evaluation job at repair depth $d$ and $R_d$ a repair job at depth $d$.
For a maximum depth $D$ and fan-out $F$, the implemented scheduler obeys these production
rules:
\[
\begin{aligned}
  E_d &\longrightarrow \text{at most one }R_{d+1}, && d<D,\\
  E_D &\longrightarrow \varnothing,\\
  R_d &\longrightarrow E_d + \text{at most }F E_{d+1}, && d<D,\\
  R_D &\longrightarrow E_D.
\end{aligned}
\]
The $E_d$ created directly by $R_d$ is verification of the repaired node. The $E_{d+1}$
jobs are re-evaluations of other nodes reached by state-change propagation. An evaluation
schedules at most one repair because findings are severity-sorted and the loop stops after
the first actionable one.

\subsection{No false clean from an incomplete run}

\begin{proposition}[Completion before closure]
\label{prop:completion}
Let $f$ be an unresolved finding associated with detector rule $r$. The verification handler
marks $f$ fixed only if (i) the run reports complete, (ii) $r$ is in the set of checked rule
identifiers, and (iii) no returned finding matches $f$. Therefore detector error or omission
of $r$ cannot by itself close $f$.
\end{proposition}

\begin{proof}
The handler constructs the required-rule set from $f$, computes missing rules by set
difference, and defines completion as reported completion conjoined with an empty missing
set. The fixed transition is guarded by that completion value and absence of a matching
finding. Failure of (i) or (ii) makes the guard false. The unresolved state is retained.
\end{proof}

This is a fail-closed statement about execution evidence. It does not establish that a
detector which ran is semantically complete. If $r$ has a false negative, condition (iii) can
hold while the defect remains.

\subsection{Finite cascade bound}

\begin{theorem}[Finite repair cascade]
\label{thm:finite}
Assume no exogenous work is introduced, every evaluation schedules at most one repair, every
repair has propagation fan-out at most $F$, and depths are restricted to
$0,1,\ldots,D$. Let $r_d$ be the maximum number of repair jobs at depth $d$ descended from a
single root $E_0$. Then
\[
 r_0=0,\qquad r_1=1,\qquad
 r_d\le r_{d-1}+F r_{d-2}\quad(2\le d\le D).
\]
The total number of evaluation and repair jobs is bounded by
\begin{equation}
  J(D,F)\le 1 + 2\sum_{d=1}^{D}r_d
                + F\sum_{d=1}^{D-1}r_d.
  \label{eq:jobbound}
\end{equation}
Hence every repair cascade is finite.
\end{theorem}

\begin{proof}
The root $E_0$ creates at most one $R_1$, so $r_1=1$. A repair at depth $d$ can arise in two
ways. First, the verification evaluation emitted by a repair at depth $d-1$ can emit one
repair at depth $d$, contributing at most $r_{d-1}$. Second, each repair at depth $d-2$
emits at most $F$ propagated evaluations at depth $d-1$, each of which can emit one repair at
depth $d$, contributing at most $F r_{d-2}$. There are no other production rules, proving
the recurrence.

Each repair job contributes itself and one verification evaluation, giving
$2\sum_{d=1}^D r_d$. Each repair below the maximum depth contributes at most $F$ propagated
evaluations, giving $F\sum_{d=1}^{D-1}r_d$. Adding the root evaluation yields
Equation~\eqref{eq:jobbound}. Every sum is finite because $D$ is finite.
\end{proof}

\begin{corollary}[Concrete LitHarness bound]
For $D=3$ and $F=12$, the recurrence gives
\[
 (r_1,r_2,r_3)=(1,1,13)
\]
and
\[
 J(3,12)\le 1+2(1+1+13)+12(1+1)=55.
\]
Thus one root evaluation can create at most 15 repair jobs and 40 evaluation jobs. Content-
derived job IDs, persistent-finding suppression, absent findings, and smaller observed
fan-out can only reduce this count.
\end{corollary}

The bound is intentionally about scheduling, not model invocations. Individual jobs also
have bounded attempt budgets, so a persistent malformed response ends in a parked or poisoned
unit instead of extending the cascade indefinitely.

\subsection{Propagation abstention}

Propagation is a routing claim, not a defect claim. A recognized fact change may reach later
scenes that mention the subject and predicate; an unrecognized change kind must return
``unhandled,'' not an empty affected set. Formally, let
\[
 \Gamma:\mathcal{C}\rightharpoonup 2^{\mathcal{V}}
\]
be a partial propagation function over change descriptions. The undefined value $\bot$ is
semantically different from $\Gamma(c)=\varnothing$. Collapsing $\bot$ to the empty set turns
``no rule applies'' into ``a rule proved no impact,'' a false completeness claim.

This distinction is valuable beyond narrative systems. Any autonomous dependency analyzer
should expose abstention as first-class state, because empty outputs are otherwise
indistinguishable from missing analysis.

\section{Risk-Controlled Craft Gates}
\label{sec:statistics}

\subsection{Craft gating as selective prediction}

Selective classification trades coverage for risk by allowing a system to reject or abstain
\cite{elyaniv2010selective}. A craft gate has the inverse operational phrasing but the same
mathematics. For unit $X$ and human defect label $Y\in\{0,1\}$, let a deterministic metric
$s(X)$ and threshold $t$ produce a flag
\[
 A_t(X)=\ind\{s(X)>t\}
\]
for an upper-tail defect. The selective precision is
\[
 p_t=\Prob(Y=1\mid A_t(X)=1).
\]
The gate's coverage, or workload, is
\[
 q_t=\Prob(A_t(X)=1).
\]

A lower bound on $p_t$ controls false refusals among flagged scenes. It says nothing about
unflagged defects unless recall is separately measured, and it says nothing about utility
unless the label itself represents the desired literary judgment.

\subsection{Evidence class and grain}

Before sampling uncertainty, a gate must match its evidence referent.

\begin{definition}[Grain compatibility]
Let $g_e$ be the grain at which the label is observed and $g_d$ the grain of the blocked
decision. Evidence is grain-compatible only if it is no coarser than the decision. A
story-level conversion label cannot license a scene-level refusal.
\end{definition}

Human judgments about the system's own scenes can support a scene-quality claim. A published
corpus percentile can support only an out-of-distribution claim. Story-level reader behavior
can rank stories but cannot identify which scene caused the behavior. This is the ecological
fallacy expressed as a type check.

\subsection{The point-estimate gap}

The current judgment-evidence rule requires:

\begin{enumerate}
  \item point precision $\widehat p\ge 0.80$;
  \item at least 50 held-out judgments; and
  \item at least 17 flagged judgments.
\end{enumerate}

The value 17 has a valid but narrow derivation. If all 17 flags are correct, the lower end of
a two-sided 95\% Clopper-Pearson interval is
\[
 (0.025)^{1/17}=0.8049>0.80.
\]
It does not follow that any observed precision above 0.80 with 17 flags has a lower bound
above 0.80.

Let $K\sim\operatorname{Binomial}(n,p)$ be the number of correct flags among $n$ flagged
holdout units. The exact two-sided $1-\alpha$ lower confidence limit is
\begin{equation}
 L_{\alpha}(K,n)=
 \begin{cases}
 0, & K=0,\\
 \BetaInv(\alpha/2;K,n-K+1), & K>0,
 \end{cases}
 \label{eq:cp}
\end{equation}
where $\BetaInv(u;a,b)$ is the $u$ quantile of a Beta$(a,b)$ distribution
\cite{clopper1934}.

\begin{table}[ht]
\centering
\caption{Point precision does not imply a 0.80 lower confidence bound. Exact limits are
two-sided 95\% Clopper-Pearson values.}
\label{tab:cp}
\begin{tabular}{rrrrl}
\toprule
Correct $K$ & Flagged $n$ & $\widehat p$ & Exact lower limit & Current rule \\
\midrule
17 & 17 & 1.000 & 0.805 & promotable \\
14 & 17 & 0.824 & 0.566 & promotable \\
40 & 50 & 0.800 & 0.663 & promotable \\
45 & 50 & 0.900 & 0.782 & promotable \\
48 & 50 & 0.960 & 0.863 & promotable \\
\bottomrule
\end{tabular}
\end{table}

The second row is an executable counterexample in the current domain model: with a holdout
size of 50, a flagged count of 17, and precision $14/17$,
\code{why\_not\_promotable} returns \code{None}. The problem is
also representational. The calibration schema stores a floating-point precision and the
number flagged, but not the integer $K$ of true-positive flags. The exact interval cannot be
reconstructed from the durable record, and an inconsistent non-integer product
$n\widehat p$ is not rejected.

\subsection{A certificate with an exact guarantee}

We propose that a judgment calibration store the tuple
\[
 C=(t,d,N,n,K,m,\alpha,p_0,q_{\min},\delta,\mathcal{H}),
\]
where $t$ is the threshold, $d$ its direction, $N$ total holdout units, $n$ flagged units,
$K$ correct flags, $m$ the number of candidate gates simultaneously safety-tested,
$\alpha$ the family-wise error budget, $p_0$ the required precision floor, $q_{\min}$ a
minimum operational coverage, $\delta$ an expiration or shift policy, and $\mathcal{H}$ a
content digest of labels, metric implementation, generator configuration, split assignment,
and selection protocol.

For a threshold fixed before the safety-test split, promotion requires
\begin{align}
  L_{\alpha}(K,n)&\ge p_0, \label{eq:single}\\
  n&\ge n_{\min},\\
  N&\ge N_{\min},\\
  n/N&\ge q_{\min},
\end{align}
plus evidence-class, grain, currency, and digest checks. A one-sided exact limit may replace
Equation~\eqref{eq:cp} if the protocol declares the directional claim in advance. The key is
that the confidence convention is part of $\mathcal{H}$ rather than chosen after seeing the
data.

\begin{theorem}[Single-candidate promotion guarantee]
\label{thm:singlecp}
Assume the flagged safety-test units are independent Bernoulli trials with constant selective
precision $p_t$, and threshold $t$ is fixed independently of the safety-test labels. If the
gate is promoted only when the exact lower limit satisfies
$L_{\alpha}(K,n)\ge p_0$, then
\[
 \Prob_{p_t}\bigl(\text{promote and }p_t<p_0\bigr)\le\alpha.
\]
\end{theorem}

\begin{proof}
Clopper-Pearson inversion gives coverage at least $1-\alpha$ for every binomial parameter:
$\Prob_{p_t}(L_{\alpha}(K,n)\le p_t)\ge1-\alpha$. If $p_t<p_0$ and the rule promotes, then
$L_{\alpha}(K,n)\ge p_0>p_t$, an event contained in the confidence-coverage failure event.
Its probability is at most $\alpha$.
\end{proof}

This is a high-probability statement over the sampled calibration data, in the spirit of
behavioral constraints on a selected algorithm \cite{thomas2019seldonian}. It is not a
distribution-free guarantee under arbitrary deployment shift.

\subsection{Threshold search and multiplicity}

If one examines many thresholds on the same holdout and records only the best, the fixed-
candidate assumption in Theorem~\ref{thm:singlecp} is false. A digest of the final verdict
set does not reveal how many candidates were tried.

The preferred protocol uses three disjoint roles:

\begin{enumerate}
  \item \textbf{development data} for metric design;
  \item \textbf{selection data} for choosing one threshold and direction; and
  \item \textbf{safety-test data} for the one final promotion test.
\end{enumerate}

When $m$ candidates must be tested on the same safety set, a conservative correction uses
$\alpha/m$ for each candidate.

\begin{theorem}[Family-wise promotion guarantee]
\label{thm:family}
Let $m$ candidate gates be tested, and suppose each candidate $j$ is promoted only if its
exact lower limit at error level $\alpha/m$ exceeds $p_0$. Under each candidate's binomial
sampling assumptions,
\[
 \Prob\bigl(\exists j:\text{ gate }j\text{ is promoted and }p_j<p_0\bigr)\le\alpha.
\]
No independence between candidates is required for this bound.
\end{theorem}

\begin{proof}
For candidate $j$, Theorem~\ref{thm:singlecp} bounds erroneous promotion by $\alpha/m$.
The union bound over $m$ candidates gives at most
$\sum_{j=1}^{m}\alpha/m=\alpha$.
\end{proof}

For a one-sided 95\% lower bound and perfect observed precision, the minimum number of flags
needed to establish $p_0=0.80$ is
\[
 n\ge\left\lceil\frac{\log(\alpha/m)}{\log(0.80)}\right\rceil.
\]
This gives 14 flags for $m=1$, 24 for $m=10$, and 35 for $m=100$. Search multiplicity is not
a philosophical caveat; it changes the sample size even when every observed flag is correct.

Conformal risk control provides a broader framework when the loss is monotone in a tuning
parameter and the desired risk is not simply binomial precision \cite{angelopoulos2024crc}.
The exact-binomial certificate is attractive here because the claim is conditional precision,
the sufficient statistics are small integers, and the resulting refusal is readily auditable.

\subsection{Exchangeability, clustering, and shift}

Scenes from the same book are not independent. They share author, generator, prompt profile,
arc position, and revision history. Treating 50 scenes from one book as 50 independent units
can make a narrow confidence interval out of one effective cluster.

The certificate should therefore declare a sampling unit and cluster identifier. Practical
options are:

\begin{itemize}
  \item split by book or generation run, never randomly by scene across the same book;
  \item report the number of independent books and generators represented;
  \item use cluster bootstrap or beta-binomial sensitivity analysis for exploratory work;
  \item reserve exact binomial promotion for a preregistered design whose Bernoulli model is
  defensible; and
  \item expire the certificate when generator, prompt profile, metric implementation, genre
  stratum, or judgment protocol changes materially.
\end{itemize}

The existing content digest and expiration date are good foundations. The proposed protocol
digest extends them to the selection procedure and configuration that generated the prose.
Without this, evidence can remain byte-current while its data-generating process has changed.

\section{The Optimization Firewall}

\subsection{Why a calibrated proxy should not become a reward}

Goodhart effects arise when selection pressure is applied to a proxy and the selected region
no longer preserves the proxy-target relationship \cite{manheim2018goodhart}. A craft gate
can create that pressure even without gradient descent. Repeatedly regenerate until the
metric passes, and the system is performing rejection sampling against the metric.

\begin{lemma}[Retry amplification]
Let independent candidate generations pass a proxy gate with probability $q\in(0,1)$. With
at most $B$ attempts, the probability that the system returns some passing candidate is
\[
 1-(1-q)^B.
\]
Conditional on success, the returned candidate distribution is supported entirely on the
proxy-passing region. This statement makes no claim about true quality.
\end{lemma}

\begin{proof}
All $B$ attempts fail with probability $(1-q)^B$; the complement is the first expression.
The stopping rule returns only after observing a pass, so every returned candidate lies in
the passing event by construction.
\end{proof}

Even $B=3$ materially increases proxy selection. If $q=0.5$, the probability of finding a
pass rises from $0.5$ to $0.875$. A weak metric is therefore not merely observed three times;
it becomes an optimizer that selects the half of the generator distribution it favors.

\subsection{Outcome classes}

The workflow should distinguish three cases.

\begin{center}
\begin{tabular}{p{0.22\textwidth}p{0.31\textwidth}p{0.35\textwidth}}
\toprule
Failure class & Permitted response & Reason \\
\midrule
Deterministic, candidate-local defect & bounded retry or located repair & The veto names a
mechanical property whose satisfaction is independently checkable. \\
Stale or malformed base & regenerate from current inputs & Repeating the same candidate
cannot resolve temporal inconsistency. \\
Calibrated craft or distribution signal & park and request judgment & Retrying would optimize
the proxy and invalidate the passive calibration claim. \\
Unknown or anonymous veto & escalate & The system lacks a justified automated action. \\
\bottomrule
\end{tabular}
\end{center}

This policy is a form of optimization firewall. The craft signal may stop a scene, but its
value does not enter the generator prompt, candidate ranking, or retry loop. A human can
inspect the parked unit, amend policy, or authorize a new plan. The proxy remains an
instrument rather than becoming the objective.

\subsection{Self-report and critic independence}

A generating model's assertion that its own output is correct is not an independent test.
The policy layer therefore rejects a blocking gate sourced from model self-report and
requires calibration identity for blocking craft gates. This structural rule is stronger
than a prompt instruction asking the model to be impartial. It removes the untrusted verdict
from the type of evidence that can authorize acceptance.

Model-family separation alone is not sufficient. Two models may share training data,
stylistic priors, and benchmark artifacts. The proposed statistical certificate therefore
tests verdicts against human labels and treats critic identity as part of the protocol digest.
Order-swapped evaluation and abstention rates are useful diagnostics, given observed judge
position bias \cite{zheng2023judge}, but diagnostics do not replace held-out human evidence.

\section{Artifact Mapping and Validation}
\label{sec:artifact}

This paper is grounded in LitHarness source commit
\code{04cf06d5af09}. Table~\ref{tab:mapping} maps each formal
claim to its implementation surface.

\begin{table}[ht]
\centering
\caption{Formal claims and repository evidence. Paths are relative to the repository root.}
\label{tab:mapping}
\small
\begin{tabular}{p{0.25\textwidth}p{0.31\textwidth}p{0.34\textwidth}}
\toprule
Claim & Implementation & Executable evidence \\
\midrule
Localized noninterference &
\code{domain/patch.py}: \code{apply\_patch}, \code{\_splice},
\code{\_verify\_preservation} &
Property-based preservation tests and patch veto tests in \code{tests/test\_domain.py} and
the repair workflow suite. \\

Stale-write rejection &
\code{domain/patch.py}: base version, base content, and span-hash checks &
Stale-base and stale-span cases in domain, draft, planner, and repair tests. \\

Completion before closure &
\code{application/repair.py}: \code{make\_evaluation\_handler} &
The incomplete-verification workflow test in
\path{tests/test_repair_workflow.py}. \\

Finite cascade &
\code{application/repair.py}: depth cap 3 and propagation fan-out cap 12 &
Depth, persistent-finding, propagation, and end-to-end repair tests in
\code{tests/test\_repair\_workflow.py}. \\

Atomic auditability &
\code{adapters/sqlite\_store.py}: \code{commit\_revision} in one immediate transaction &
Store, event, state-retraction, and conductor crash/replay tests. \\

Promotion provenance &
\code{domain/calibration.py}, migrations 014, 015, and 018 &
69 calibration tests covering evidence class, grain, expiration, digests, controls, and
promotion refusal. \\

Optimization firewall &
\code{domain/policy.py}: \code{PARKABLE}, \code{decide}, and the policy-decision
constructor &
Mixed-veto, park, retry, self-report, and uncalibrated-critic tests. \\
\bottomrule
\end{tabular}
\end{table}

On 17 August 2026, the repository collected 795 tests. The full default suite completed with
787 passed and 8 skipped in 86.52 seconds on the project machine. This establishes that the
mapped executable properties are mutually compatible in that environment. It does not turn
unit tests into a proof of the Python runtime, SQLite, cryptographic primitives, or detector
semantics. The mathematical proofs are conditional on the stated implementation contracts;
the tests check that the current artifact follows those contracts across a broad set of
examples and property-generated inputs.

\subsection{Executable statistical counterexample}

The following construction uses only the public calibration domain object:

\begin{verbatim}
c = Calibration(
    calibration_id="cal-demo", metric_id="m",
    holdout_size=50, precision=14/17,
    threshold=1.0, direction=Direction.ABOVE,
    verdicts_digest="d", measured_at="2026-08-17",
    evidence_class=EvidenceClass.JUDGMENT,
    grain=Grain.UNIT, flagged=17,
)
assert c.why_not_promotable(
    "2026-08-17", "d", answered=50
) is None
\end{verbatim}

The assertion passes. This does not contradict the source comments; they explicitly state
that the 17-flag floor does not make 0.80 a confidence bound in general. The contribution
here is to turn that caveat into a complete promotion rule and identify the schema fields
needed to enforce it.

\section{Recommended Minimal Change Set}

No quality metric should be added merely to exercise this design. The following changes are
useful when the first real human-labeled calibration exists.

\subsection{Calibration schema}

Add the fields below, with old rows defaulting to an explicitly unpromotable protocol class.

\begin{itemize}
  \item \code{true\_positives INTEGER}: the integer $K$.
  \item An integer false-positive count: redundant with flagged count but valuable for
  direct integrity checks; require $K+FP=n$.
  \item \code{confidence\_alpha REAL} and \code{confidence\_sidedness TEXT}.
  \item \code{precision\_lower\_bound REAL}: derived and verified on read, not trusted as an
  independent input.
  \item \code{selection\_family\_size INTEGER} and \code{selection\_protocol\_digest TEXT}.
  \item \code{split\_digest TEXT}: content address over unit-to-split assignment.
  \item \code{generator\_config\_digest TEXT}, \code{critic\_config\_digest TEXT}, and a
  declared cluster key such as book or generation run.
\end{itemize}

Do not infer $K$ by rounding \code{precision * flagged}; doing so turns a malformed record
into apparently exact evidence. Compute and store the point precision from counts instead.

\subsection{Promotion predicate}

For judgment evidence, replace the point-estimate check with:

\begin{verbatim}
require integer confusion counts are internally consistent
require evidence grain covers decision grain
require threshold was fixed outside the safety-test labels
alpha_each = alpha_family / selection_family_size
lower = exact_binomial_lower(K, flagged, alpha_each, sidedness)
require lower >= required_precision
require holdout_size >= minimum_holdout
require flagged / holdout_size >= minimum_coverage
require evidence and configuration digests are current
\end{verbatim}

The operation should return the exact reason for refusal, as the current promotion path
already does. A failed statistical certificate should remain advisory and visible; it should
not silently disappear.

\subsection{Testing strategy}

Add property tests for:

\begin{enumerate}
  \item monotonicity of the lower bound in $K$ for fixed $n$;
  \item non-promotion of every $(K,n)$ pair whose exact lower bound is below $p_0$;
  \item refusal of non-integer or inconsistent confusion counts;
  \item multiplicity making promotion weakly harder as $m$ increases;
  \item byte changes in split assignment, generator configuration, or selection protocol
  changing the certificate identity; and
  \item backward compatibility making pre-migration rows advisory rather than inventing
  missing evidence.
\end{enumerate}

\section{Limitations and Falsification Criteria}

\subsection{Limits of the deterministic proofs}

The patch theorem is only as strong as the target intervals. A detector can cite the wrong
passage, and the repair can perfectly preserve everything outside that wrong passage. Hashes
prove identity, not relevance. Likewise, a propagation rule can route evaluation to the wrong
set of nodes. The existing perfect propagation score is on four in-sample development cases
and is not evidence of generalization \cite{litharnessbrief}.

The cascade bound excludes exogenous work such as a new director directive, a newly accepted
draft, or manual revival. These are new roots, not descendants of the bounded root. The bound
also assumes the scheduler respects depth fields and fan-out caps; arbitrary database
tampering is outside the threat model.

\subsection{Limits of the statistical guarantee}

The exact interval controls binomial sampling uncertainty, not label validity, dependence,
or distribution shift. The result is falsified as a deployment claim if any of the following
occurs:

\begin{itemize}
  \item blocked units are at a different grain from the labels;
  \item threshold selection examined the safety-test labels without correction;
  \item scenes are treated as independent although a few books dominate the sample;
  \item generator, prompt, genre, or judgment protocol changes without invalidating the
  certificate;
  \item a human audit on a preregistered fresh sample yields a lower bound below $p_0$; or
  \item the metric is fed back into generation, changing the deployment distribution from
  the passive one on which it was calibrated.
\end{itemize}

These are not post hoc caveats. Each should be represented in the certificate schema or in a
test that makes promotion fail.

\subsection{Literary validity remains open}

The most important qualities in long-form fiction are causal and longitudinal: a scene must
change something, progression must cost something, promises must be paid, and voice must
remain particular. A confidence interval around a shallow metric does not make the label more
meaningful. The correct outcome when no valid metric exists is an empty blocking-craft table,
not a weaker promotion standard.

The research direction that survives is intervention plus human attribution: construct
minimal controlled edits that alter one dramatic property, collect blinded judgments at the
same grain, and test whether a proposed signal distinguishes the intervention from a
placebo. Until such labels exist, the statistical machinery in this paper should protect the
empty state rather than pressure the project to fill it.

\section{Discussion}

The central design choice is to certify the boundary around creativity instead of claiming
to certify creativity itself. This is analogous to runtime assurance for an advanced but
unverified controller: the high-performance component remains free to propose actions, while
a smaller trusted layer restricts the state transitions those actions may cause. Unlike a
physical controller, a manuscript has no complete safety specification. The monitor therefore
enforces only properties that are both meaningful and decidable: scope, freshness,
attribution, completion, and boundedness.

This separation produces a useful asymmetry. Deterministic failures may drive repair because
their postcondition is independently testable. Statistical craft failures park because a
retry would change the population the statistic describes. The same gate can be valid as a
passive alarm and invalid as an optimization target.

The finite-sample result also illustrates a broader engineering rule. Evidence provenance is
not only a digest of the observations. It includes how a candidate was selected. Recording
50 judgments and their hash does not reveal whether one threshold was specified in advance or
100 thresholds were tried and the best retained. A reconstructible safety claim needs counts,
splits, candidate-family size, configurations, grain, and time.

Finally, a bounded autonomous system need not pretend that every refusal is recoverable by
more autonomy. Park and abstain are legitimate outcomes. Selective-prediction theory treats
coverage as a first-class trade-off; an authoring system should do the same. The ability to
say ``this unit requires judgment'' is not a failure of automation. It is the control that
makes the rest of the automation defensible.

\section{Conclusion}

Autonomous long-form generation cannot presently prove that its prose is good. It can prove
substantial facts about how uncertain prose becomes part of a book. Interval-certified
patches preserve uncited text; content addresses reject stale work; transactional acceptance
keeps revisions, decisions, state, events, and verification jobs together; incomplete
detector runs cannot close findings; and explicit depth and fan-out limits make repair
cascades finite. In the current LitHarness configuration, the worst-case descendant count
from one evaluation is 55 jobs.

Craft gates require a different proof obligation. Point precision and a fixed minimum number
of flags are not a confidence guarantee. Exact count-based lower bounds, selection-test
separation, multiplicity correction, grain compatibility, and configuration provenance turn
promotion from a number in a row into a reconstructible statistical certificate. The gate
must then remain outside the optimization loop.

The resulting architecture makes a modest claim with strong consequences: an untrusted model
may propose, but only a small deterministic and statistically accountable kernel may decide
what the proposal is allowed to change and what evidence is sufficient to stop the book.

\appendix

\section{Closed Form and Growth of the Cascade Recurrence}

The recurrence $r_d=r_{d-1}+F r_{d-2}$ has characteristic equation
\[
 x^2=x+F,
\]
with roots
\[
 \lambda_{\pm}=\frac{1\pm\sqrt{1+4F}}{2}.
\]
For the equality case with $r_0=0$ and $r_1=1$,
\[
 r_d=\frac{\lambda_+^d-\lambda_-^d}{\lambda_+-\lambda_-}.
\]
Thus work grows exponentially in $D$ for fixed $F>0$, even though it remains finite. This is
why both controls are necessary: a depth limit alone permits a large finite burst when $F$
is large, and a fan-out limit alone permits an infinite chain when $D$ is absent.

For $F=12$, $\lambda_+=4$ and $\lambda_-=-3$, so
\[
 r_d=\frac{4^d-(-3)^d}{7}.
\]
At $D=3$ this yields $(1,1,13)$ as above. Raising the depth to 4 would yield $r_4=25$ and
increase the total bound from 55 to
\[
 1+2(1+1+13+25)+12(1+1+13)=261.
\]
A one-step depth change therefore has a sixfold operational effect in this configuration.

\section{Exact-Bound Reference Values}

Table~\ref{tab:reference} gives reference values useful for implementation tests. Values use
the two-sided 95\% Clopper-Pearson lower limit from Equation~\eqref{eq:cp}.

\begin{table}[ht]
\centering
\caption{Reference exact lower limits for implementation tests.}
\label{tab:reference}
\begin{tabular}{rrrr}
\toprule
$K$ & $n$ & $K/n$ & $L_{0.05}(K,n)$ \\
\midrule
17 & 17 & 1.00 & 0.804936 \\
14 & 17 & 0.823529 & 0.565682 \\
16 & 20 & 0.80 & 0.563386 \\
40 & 50 & 0.80 & 0.662817 \\
45 & 50 & 0.90 & 0.781865 \\
48 & 50 & 0.96 & 0.862862 \\
49 & 50 & 0.98 & 0.893530 \\
50 & 50 & 1.00 & 0.928878 \\
90 & 100 & 0.90 & 0.823777 \\
\bottomrule
\end{tabular}
\end{table}

An observed precision exactly equal to $p_0$ cannot have a nontrivial lower confidence bound
at least $p_0$ for finite $n$. Consequently, a rule that requires both
$\widehat p\ge p_0$ and $L\ge p_0$ will promote only observations strictly above $p_0$, apart
from degenerate confidence levels. This is expected: uncertainty requires margin.

\section{Safety-Case Checklist}

For a new automated repair or gate, the following questions are sufficient to expose most
category errors discussed in this paper.

\begin{enumerate}
  \item What exact bytes or nodes may change, and is the untouched complement checked?
  \item What binds the action to the version it was generated against?
  \item Can detector failure be represented separately from a clean result?
  \item What is the maximum descendant work from one accepted change?
  \item Does an unsupported case abstain, or collapse to an empty answer?
  \item What is the label's grain, and what is the decision's grain?
  \item Are integer confusion counts stored, or only rounded rates?
  \item Was the tested threshold fixed without access to the safety labels?
  \item How many candidate gates were examined?
  \item Which generator, critic, prompt, stratum, and split produced the evidence?
  \item Does failure cause observation, parking, repair, regeneration, or optimization?
  \item What event or digest makes every one of these answers reconstructible later?
\end{enumerate}

\begin{thebibliography}{99}

\bibitem{angelopoulos2024crc}
A. N. Angelopoulos, S. Bates, A. Fisch, L. Lei, and T. Schuster.
Conformal risk control.
\emph{International Conference on Learning Representations}, 2024.
\url{https://arxiv.org/abs/2208.02814}

\bibitem{clopper1934}
C. J. Clopper and E. S. Pearson.
The use of confidence or fiducial limits illustrated in the case of the binomial.
\emph{Biometrika}, 26(4):404--413, 1934.
\url{https://doi.org/10.1093/biomet/26.4.404}

\bibitem{elyaniv2010selective}
R. El-Yaniv and Y. Wiener.
On the foundations of noise-free selective classification.
\emph{Journal of Machine Learning Research}, 11:1605--1641, 2010.
\url{https://jmlr.csail.mit.edu/papers/v11/el-yaniv10a.html}

\bibitem{huang2024selfcorrect}
J. Huang, X. Chen, S. Mishra, H. S. Zheng, A. W. Yu, X. Song, and D. Zhou.
Large language models cannot self-correct reasoning yet.
\emph{International Conference on Learning Representations}, 2024.
\url{https://arxiv.org/abs/2310.01798}

\bibitem{kung1981optimistic}
H. T. Kung and J. T. Robinson.
On optimistic methods for concurrency control.
\emph{ACM Transactions on Database Systems}, 6(2):213--226, 1981.
\url{https://doi.org/10.1145/319566.319567}

\bibitem{ligatti2005edit}
J. Ligatti, L. Bauer, and D. Walker.
Edit automata: Enforcement mechanisms for run-time security policies.
\emph{International Journal of Information Security}, 4(1--2):2--16, 2005.
\url{https://doi.org/10.1007/s10207-004-0046-8}

\bibitem{manheim2018goodhart}
D. Manheim and S. Garrabrant.
Categorizing variants of Goodhart's law.
\emph{arXiv:1803.04585}, 2018.
\url{https://arxiv.org/abs/1803.04585}

\bibitem{thomas2019seldonian}
P. S. Thomas, B. Castro da Silva, A. G. Barto, S. Giguere, Y. Brun, and E. Brunskill.
Preventing undesirable behavior of intelligent machines.
\emph{Science}, 366(6468):999--1004, 2019.
\url{https://doi.org/10.1126/science.aag3311}

\bibitem{zheng2023judge}
L. Zheng, W.-L. Chiang, Y. Sheng, S. Zhuang, Z. Wu, Y. Zhuang, Z. Lin, Z. Li,
D. Li, E. P. Xing, H. Zhang, J. E. Gonzalez, and I. Stoica.
Judging LLM-as-a-judge with MT-Bench and Chatbot Arena.
\emph{Advances in Neural Information Processing Systems}, 36, 2023.
\url{https://arxiv.org/abs/2306.05685}

\bibitem{litharnessbrief}
LitHarness contributors.
The unsolved quality-measurement problem in LitHarness.
Internal research brief, revision of 17 August 2026,
\code{research/quality-measurement/BRIEF.md}.

\bibitem{litharnessartifact}
LitHarness contributors.
LitHarness autonomous book-production system, source artifact.
Commit \code{04cf06d5af09eb1c9b11e9a59968ad179912e4cf}, 17 August 2026.

\end{thebibliography}

\end{document}
